\section{Question, result, and mathematical ownership}
\label{sec:scope}

The question is whether a proposed positive physical sector can be obtained
from an action with an indefinite ambient pairing. The answer depends on
its real fields, stationary background, constraints, boundary data,
observation map, and Hamiltonian domain. A gauge-group designation or a
positive auxiliary norm does not determine these objects.

The principal result of this manuscript is action-specific. We fix the
connection metric on $Y=\operatorname{Met}_{1,3}(X)$ and the
commutator/two-symmetrized-channel realization of the first bosonic action
$I_1^B$ in Weinstein's draft \cite{GU}. The coordinate metric, signature
convention, bracket normalization, and boundary polarization below are
explicit reconstruction choices. They are not claimed to be uniquely
selected by that source. The source's printed curvature residual, its
square $I_2^B$, the observing metric, and the metric-fibre coordinate remain
separate objects.

\begin{theorem}[Main result; precise constructions below]
\label{thm:main}
For the action \eqref{eq:action}, the stationary branches specified in
\cref{sec:background} have the following properties.
\begin{enumerate}
\item Their imaginary coefficient-conjugation packet has a smooth,
compactly supported constraint-reduced phase core on which its quadratic
Noether Hamiltonian takes values tending to both infinities.
\item Their real packet has such families with zero metric perturbation
and zero full linearized metric forcing. On a spatial torus and the
reflection-invariant fibre patch of \cref{sec:real}, these families satisfy
all finite-order formal compatibility conditions of the odd-Dirichlet
polarization in \cref{sec:boundary}.
\item Every completion containing the respective core and agreeing with
its quadratic energy form inherits its two-sided unboundedness. Bare
smooth section-germ observation does not make the imaginary-packet
Hamiltonian descend to an observed quotient.
\end{enumerate}
\end{theorem}

The proofs are given in \cref{sec:imaginary,sec:real,sec:selection}.
The theorem neither presupposes nor constructs a global Cauchy evolution.
In particular, compatible formal time jets need not converge. No
Hilbert-space normalization, quantum spectral theorem, nonlinear runaway,
or integrability of every linearized phase datum into nonlinear physical
data is inferred. The two coupling branches are two selected values with
opposite signs, not all prescribed couplings or all backgrounds.

\begin{center}\small
\begin{tabular}{p{0.21\linewidth}p{0.71\linewidth}}
\toprule
Object & Convention and role \\
\midrule
Carrier & Clifford-valued one-forms on $Y$, a base metric $g(x)$, and
source gauge variables removed by exact dressing. \\
Real structure & A phased anti-involution slice; its real and imaginary
coefficient packets are defined in \eqref{eq:packets}. \\
Pairings & Scalar Clifford trace and exterior top-degree action pairing;
reduced time symplectic form and Noether energy are distinguished. \\
Grading & Clifford grade, coefficient conjugation, primary time kernel,
and differential order; no coefficient rank is a particle count. \\
Domain & Local product patches with stated variations; compact phase
cores and one explicit polarization, without a closed physical generator. \\
Target & Semiboundedness of the selected quadratic energy and its
observation descent, not every source mechanism. \\
\bottomrule
\end{tabular}
\end{center}

The distinction between a mathematical reconstruction and a physical GU
realization remains throughout. Ordinary Yang--Mills and external scalar
models in \cref{sec:controls} bind their own actions. Transporting one of
their conclusions would require a map of carriers and real structures,
intertwining the actions, constraints, domains, and observable in question.
The separate quantum/PDE controls of \cref{sec:quantum} provide no such
identification. No literature-priority claim is made for this manuscript.

\section{Canonical geometry and an exact stationary background}
\label{sec:background}

\subsection{The observing metric and the fibre metric}
Let $X$ be a four-dimensional coordinate patch and let $h$ denote the
Lorentzian symmetric matrix at a point of its metric bundle. The independent
observing metric is $g$. Set
\begin{align}
 G_Y(g)&=h_{ij}\dd x^i\dd x^j+D_h(\theta,\theta),\label{eq:geometry}\\
 D_h(A,B)&=\tr(h^{-1}Ah^{-1}B)
       -\tfrac12\tr(h^{-1}A)\tr(h^{-1}B),\nonumber\\
 \theta_{ab}&=\dd h_{ab}
 -(\Gamma(g)^c{}_{ia}h_{cb}+\Gamma(g)^c{}_{ib}h_{ac})\dd x^i.
 \nonumber
\end{align}
We use $g_*=\eta=\diag(1,-1,-1,-1)$ and the resulting $(7,7)$ ambient
signature. This plus-first convention is fixed for the calculation; a
source-convention bridge is not established merely by writing the same
signature name. The coordinate density in symmetric fibre coordinates is
\begin{equation}
 \dd\mu_Y=\mu(h)\dd^4x\dd^{10}h,
 \qquad \mu(h)=8|\det h|^{-2}.\label{eq:density}
\end{equation}
It is independent of $g$ at fixed $h$.

Although $\Gamma(g_*)=0$, the coefficients depend on $h$.
The Levi-Civita connection obtained from \eqref{eq:geometry} has
\begin{equation}
 \Ric_Y^\sharp=\tfrac14(P_H+P_r)-\tfrac54P_0,
 \qquad \Scal(G_Y)=-10.\label{eq:ricci}
\end{equation}
Here $P_H$ is horizontal, $P_0$ is vertical traceless, and $P_r$ is the
vertical trace projector. Their ranks are $4,9,1$. The vertical field
$r=(0,h/2)$ satisfies
\begin{equation}
 G_Y(r,r)=-1,\qquad P_rz=-G_Y(r,z)r,
 \qquad \nabla r=\tfrac14P_H.\label{eq:radial}
\end{equation}
The coordinate Koszul calculation checks torsion, metricity, Bianchi, and
these contractions \cite{GeometryAssessment}. For example, a vertical
traceless plane has sectional curvature $-1/2$. Thus setting the spin
curvature to zero would change the geometry.

Constant rescaling of $g$ leaves $\Gamma(g)$, and hence $G_Y$ at fixed
$(x,h)$, unchanged. Moving $h\mapsto sh$ instead gives the familiar
$-8$ logarithmic density derivative in fixed coordinates. That fibre
motion is not an observing-metric Euler variation. This distinction
corrects the interpretation of the flat controls in \cref{sec:flat}.

\subsection{The selected action and its real slice}
Let $B$ be the spin Levi-Civita connection of $G_Y$, $T$ its dressed
connection displacement, and $\gamma$ Clifford multiplication with
$\gamma_a\gamma_b+\gamma_b\gamma_a=2\eta_a\delta_{ab}$ in an
orthonormal frame. The selected functional is
\begin{equation}
 I=\int_Y\Tr\left[
 T\wedge\cS(F_B)+\tfrac12T\wedge\cS(D_BT)
 +\tfrac13T\wedge\cS(T\wedge T)
 +\tfrac\kappa2T\wedge *T\right].\label{eq:action}
\end{equation}
The exterior integral uses scalar Clifford trace, normalized by
$\tr(1)=1$. The following contraction specifies the Shiab normalization
completely by linearity. For $p<q$ and Clifford coefficients $C,D$,
\begin{align}
 &\frac{\Tr[e^\mu C\wedge\cS(e^p\wedge e^qD)]}{\operatorname{vol}_Y}
   =\eta_p\eta_q\left(\delta_{\mu p}\tr(C[\gamma_q,D])
       -\delta_{\mu q}\tr(C[\gamma_p,D])\right.\nonumber\\[-0.3em]
 &\hspace{7em}\left.
       -\frac{\eta_\mu}{2}\tr(C\{\gamma_\mu,\{\gamma_p\gamma_q,D\}\})
       \right).\label{eq:shiab}
\end{align}
The last sign includes $*\operatorname{vol}_Y=-1$ and the two factors
of $i$ in the symmetrized channels. Equation \eqref{eq:action} is
motivated by source Eq.~(9.4), p.~44, but the source's displayed residual
is not substituted for its Euler derivative.

Define a conjugate-linear Clifford anti-involution $\mathfrak j$ by
$\mathfrak j(\gamma_a)=-\gamma_a$ and
$\mathfrak j(AB)=\mathfrak j(B)\mathfrak j(A)$. The source-phased real
slice is $\mathcal U=\{Z:\mathfrak j(Z)=-Z\}$. Under ordinary coefficient
conjugation it splits as
\begin{align}
 R&=\Lambda^1+\Lambda^2+\Lambda^5+\Lambda^6
       +\Lambda^9+\Lambda^{10}+\Lambda^{13}+\Lambda^{14},\nonumber\\
 J&=i(\Lambda^0+\Lambda^3+\Lambda^4+\Lambda^7
                  +\Lambda^8+\Lambda^{11}+\Lambda^{12}).\label{eq:packets}
\end{align}
All grade spaces are real before the displayed $i$. Commutators and
$i$ times anticommutators preserve $\mathcal U$; Hodge and the spin
connection preserve its exterior extension. For $A,B\in\mathcal U$,
$\overline{\tr(AB)}=\tr(\mathfrak j(AB))=\tr(BA)=\tr(AB)$.
The action is real on this slice. This involution is not a positive
Hilbert adjoint.

\subsection{Stationarity before restriction}
Consider the intrinsic field
\begin{equation}
 T_*(z)=a\gamma((P_H+P_0)z)+c\gamma(P_rz)
 -\frac v8\gamma(r)\gamma(P_Hz)+w\gamma(r)\gamma(P_0z).
 \label{eq:background}
\end{equation}
The last two terms have grade two. Vary \eqref{eq:action} in arbitrary
one-form directions before inserting \eqref{eq:background}.
For a test variation $U$, the derivative term contributes
\begin{equation}
 \frac12\sum_s\left\{
 \langle U,\cS(e^s\wedge\nabla_sT)\rangle
 -\langle\nabla_sT,\cS(e^s\wedge U)\rangle\right\};
 \label{eq:formal-adjoint}
\end{equation}
the second summand is essential. The cubic derivative has both
$\langle U,\cS(T\wedge T)\rangle/3$ and
$\langle T,\cS(U\wedge T+T\wedge U)\rangle/3$.
The complete distortion covector is zero exactly when the following
five representative rows vanish:
\begin{align}
16E_H={}&4224a^2+768ac+16a\kappa+v^2-72vw-3v
                   +768w^2+108w-84,\nonumber\\
24E_0={}&6336a^2+1152ac+24a\kappa+3v^2-128vw-12v
                   +896w^2+168w-90,\nonumber\\
24E_r={}&7488a^2+24c\kappa+v^2-144vw-6v
                   +864w^2+216w-126,\label{eq:stationary}\\
24E_{Hr}={}&132av-3168aw+12a+4cv-288cw-12c+3\kappa v,
 \nonumber\\
3E_{0r}={}&22av-352aw+3a+2cv-24cw-3c-3\kappa w.\nonumber
\end{align}
There are respectively $4,9,1,4,9$ possibly nonzero rows.
To see completeness, label $e^\mu\gamma_I$ by the XOR of its exterior
index mask and its Clifford mask. The background has only labels $0$
and $\{10\}$, where index $10$ denotes $r$.
The Clifford/exterior trace selection rule leaves 28 candidate receivers;
the scalar receiver $e^{10}1$ vanishes identically. Direct evaluation
checks every remaining linear, mass, and curvature receiver. Thus these
are full distortion equations, not four ansatz equations.

The stationary construction \cite{Stationary} gives the exact algebraic branch, specified without decimal root finding.
Put $a=\kappa\alpha$, $c=\kappa\beta$ and
\begin{align}
 D&=22v^2-220vw-15v-3168w^2+732w,\nonumber\\
 \alpha&=-\frac{3(2v^2-20vw-3v-288w^2-12w)}{8D},\nonumber\\
 \beta&=\frac{3(22v^2-220vw+3v-3168w^2+12w)}{8D}.
 \label{eq:alphabeta}
\end{align}
The conic and coupling equations are
\begin{align}
 3v^2-40vw-512w^2-15v+12w+72&=0,\nonumber\\
 V_HR_0-V_rH_0&=0,\qquad \kappa^2=-V_H/H_0,\label{eq:conic}\\
 H_0&=4224\alpha^2+768\alpha\beta+16\alpha,\nonumber\\
 R_0&=7488\alpha^2+24\beta,\nonumber\\
 V_H&=v^2-72vw-3v+768w^2+108w-84,\nonumber\\
 V_r&=v^2-144vw-6v+864w^2+216w-126.\nonumber
\end{align}
For an explicit univariate specification let $w_*$ be the unique root in
\begin{equation}
 \frac{227427945967}{10^{12}}<w_*<\frac{227427945969}{10^{12}}
 \label{eq:rootinterval}
\end{equation}
of $P(w)=\sum_{j=0}^{10}p_jw^j$, with coefficients listed in
\cref{app:algebra}. There $v(w)$ is also given explicitly.
Sturm counting certifies one root; rational intervals exclude every
parameter denominator from zero and give
$2.873978<\kappa_*^2<2.874152$.
All five rows of \eqref{eq:stationary} reduce exactly to zero in
$\QQ(w_*)$, with the square root handled by $a=\kappa\alpha$ and
$c=\kappa\beta$. Take $\kappa=\kappa_*>0$ or its negative and the
corresponding $a,c$. The map
\begin{equation}
 (a,c,\kappa)\longmapsto(-a,-c,-\kappa),\qquad (v,w)\longmapsto(v,w)
 \label{eq:mirror}
\end{equation}
preserves the common zero set. Decimal values on the positive branch,
for orientation only, are $(a,c,v,w,\kappa)\simeq
(0.151945,-0.195077,3.79821,0.227428,1.69531)$.

\begin{proposition}[Local source-variable stationarity]
The fields \eqref{eq:background} with the certified parameters are stationary
for the selected bulk functional on a local product patch, with compact
base-metric variations and compact interior distortion and gauge variations.
\end{proposition}
\begin{proof}
Translations of $x$ and simultaneous affine/congruence transformations of
$(x,h)$ preserve \eqref{eq:geometry} at constant $g$. Congruence is transitive
on the signature patch and transports the natural projectors and field.
Consequently the exact distortion equations transport from one fibre point
to the smooth patch, with local oriented spin lifts.

For the source gauge field $\eps$, set
$B_\eps=\eps^{-1}D_0\eps$ and
$\tau=\eps\varpi\eps^{-1}-(D_0\eps)\eps^{-1}$, where $D_0$ is the
reference spin connection. Conjugating curvature, covariant derivative,
and Shiab, and using cyclic trace, gives
$I(g,\eps,\varpi)=I(g,1,\tau)$ identically. The unrestricted distortion
Euler covector is zero, so the gauge-field first variation is zero.

Finally $g$ enters the local action only through $\Gamma(g)$ and its
jets. At constant $g$ its variation contains at least one base derivative
of $\delta g$; the coefficients depend on $h$ but not $x$.
Integration by parts in $x$ against compact base tests gives zero.
An induced variation of $T_*(g)$ contributes zero because its distortion
Euler equation is already satisfied. The fixed relatively compact fibre
patch makes the variational integrals finite. No integration over the
entire noncompact fibre or moving fibre boundary is used.
\end{proof}

The coupling here is a selected background parameter, not a physical
mass prediction. The failed smaller families remain informative: the
three-projector grade-one ansatz forces $a=b=c=t$ and then requires both
$312t^2+\kappa t=21/4$ and $312t^2+\kappa t=15/4$.
In \eqref{eq:background}, setting $w=0$ instead gives
$E_0-E_H=(v-5/2)^2/16+71/64>0$.
Neither exclusion applies to the enlarged family just constructed.
